% section: 漸近解析：一般項がなくても構造を知る

\section{漸近解析：一般項がなくても構造を知る}

\subsection{単調性と有界性}

数列 $a_n$ が
\[
  a_{n+1}\geq a_n
\]
を満たすとき、単調増加という。
また、ある定数 $M$ が存在して
\[
  a_n\leq M
\]
となるとき、上に有界という。

単調増加かつ上に有界な数列は収束する。
これは実数の連続性に関わる基本的な定理である。
証明の考え方は、数列の値が作る集合の上限 $L$ を取ることにある。
もし $a_n$ が $L$ より離れたままなら、$L$ が上限であることに反する。

漸化式の極限を調べるとき、
\[
  \text{単調性}+\text{有界性}
\]
を確認するのは、この定理を使うためである。
不動点方程式だけでは収束を証明できないが、単調性と有界性を加えると、極限の存在を確保できる場合がある。

\subsection{収束先を不動点方程式で絞る}

\[
  a_{n+1}=f(a_n)
\]
で、$a_n\to L$ だとする。
$f$ が連続なら、
\[
  L=f(L)
\]
が必要である。
まず不動点を求め、次に数列が実際にその不動点へ近づくことを示す。

例えば、
\[
  a_{n+1}=\frac{a_n+2}{2}
\]
なら、不動点は $L=2$ である。
$b_n=a_n-2$ と置くと、
\[
  b_{n+1}=\frac12b_n
\]
だから、
\[
  a_n=2+\frac{a_0-2}{2^n}
\]
となり、確かに $2$ へ収束する。

\subsection{縮小写像としての収束証明}

ある区間 $I$ 上で、
\[
  |f(x)-f(y)|\leq q|x-y|
  \qquad(0<q<1)
\]
が成り立つとする。
このような関数を縮小写像という。

反復 $a_{n+1}=f(a_n)$ に対して、
\[
  |a_{n+1}-a_n|
  \leq q|a_n-a_{n-1}|
  \leq q^n|a_1-a_0|
\]
である。
従って、後の項同士の距離は、等比的に小さくなる。

実際、$m>n$ に対して、
\[
  \begin{aligned}
    |a_m-a_n|
     & \leq
    |a_m-a_{m-1}|+\cdots+|a_{n+1}-a_n| \\
     & \leq
    (q^{m-1}+\cdots+q^n)|a_1-a_0|      \\
     & \leq
    \frac{q^n}{1-q}|a_1-a_0|.
  \end{aligned}
\]
$n$ を大きくすれば右辺は $0$ に近づくので、数列は一つの値へ近づく。
その極限を $L$ とすれば、連続性から $L=f(L)$ である。

これは、単に不動点を見つけるだけでなく、近くの点が縮む仕組みから収束を保証する考え方である。
ニュートン法の局所的な収束や、線形漸化式の $|r|<1$ も、この縮みという見方で統一できる。

\subsection{支配的な項を見つける}

フィボナッチ数列の一般項は、
\[
  F_n=\frac1{\sqrt5}\left(\frac{1+\sqrt5}{2}\right)^n
  -\frac1{\sqrt5}\left(\frac{1-\sqrt5}{2}\right)^n
\]
である。
二つの根の絶対値を比べると、
\[
  \left|\frac{1-\sqrt5}{2}\right|<1
\]
だから、後半の項は $n$ が大きくなると消えていく。
従って、
\[
  F_n\sim
  \frac1{\sqrt5}
  \left(\frac{1+\sqrt5}{2}\right)^n
\]
となる。

一般に、線形漸化式では特性根の絶対値が最大のものが、十分大きな $n$ での成長を支配する。
同じ絶対値の根が複数あれば、振動や $n$ の因子が残る。
一般項を完全に書くことと、長期的な成長を理解することは、関連しているが別の仕事である。

\subsection{漸近記号を使う}

\[
  a_n\sim b_n
\]
とは、
\[
  \lim_{n\to\infty}\frac{a_n}{b_n}=1
\]
という意味である。
また、
\[
  a_n=O(b_n)
\]
とは、ある定数 $C$ と十分大きな $n$ に対して、
\[
  |a_n|\leq C|b_n|
\]
となることをいう。

例えば、
\[
  n^2+3n+1\sim n^2
\]
である。
最高次の項が成長を支配するからである。
漸化式の解析では、正確な一般項を得る前に、どの程度の大きさかを見積もることが有効になる。
